\reading{CR-12}{Credit Risk (Hull Ch.24)}{STRIKE FIRST}{8}

\begin{crkeybox}[Learning objectives for this reading]
\begin{enumerate}[label=\textbf{\alph*.},leftmargin=2em]
\item Assess the credit risks of derivatives.
\item Define credit valuation adjustment (CVA) and debt valuation adjustment (DVA).
\item Calculate the probability of default using credit spreads.
\item Describe, compare, and contrast various credit risk mitigants and their role in credit analysis.
\item Describe the significance of estimating default correlation for credit portfolios and distinguish between reduced form and structural default correlation models.
\item Describe the Gaussian copula model for time to default and calculate the probability of default using the one-factor Gaussian copula model.
\item Describe how to estimate credit VaR using the Gaussian copula and the CreditMetrics approach.
\end{enumerate}
\end{crkeybox}

\begin{crtrapbox}[Single-layer reading]
CR-12 exists only in the Crash layer.
\end{crtrapbox}

% =====================================================================
\lo{CR-12 a}{Assess the credit risks of derivatives.}

Bilaterally cleared derivatives are governed by the ISDA Master Agreement. For
transactions between financial institutions, variation margin is exchanged and
initial margin is posted with third parties.

\begin{crkeybox}[When the non-defaulting party actually loses]
There is a potential loss to the non-defaulting party in two situations:

\begin{enumerate}
\item The total value of transactions is \term{positive to the non-defaulting
party} and \emph{exceeds} the collateral posted by the defaulting party.
\item The total value of transactions is \term{positive to the defaulting party},
and the collateral posted by the non-defaulting party is \emph{greater} than that
value.
\end{enumerate}

The second case is the one that is easy to miss: the non-defaulting party can be
out of the money on the trades and still lose, through over-posted collateral it
cannot recover.
\end{crkeybox}

% =====================================================================
\lo{CR-12 b}{Define credit valuation adjustment (CVA) and debt valuation adjustment (DVA).}

\begin{crdefbox}[The two adjustments]
\term{Credit valuation adjustment (CVA)} is the present value of the expected cost
to the bank of a default by the counterparty.

\smallskip
\term{Debt valuation adjustment (DVA)} is the present value of the expected cost
to the counterparty of a default by the bank. The bank's potential default is
\emph{advantageous to the bank}, as it implies a chance that the bank may not need
to fulfil required payments on its derivatives. DVA, a cost to the counterparty,
therefore increases the value of the derivatives portfolio to the bank.
\end{crdefbox}

\begin{crdefbox}[No-default value and credit-adjusted value]
The \term{no-default value} of outstanding transactions is their value assuming
neither side will default. Derivatives pricing models such as
Black-Scholes-Merton provide no-default values. Write it $f_{nd}$.

\smallskip
The \term{credit-adjusted value}, that is, the value of the outstanding
transactions when possible defaults are taken into account, is
\[ f_{nd} - \mathrm{CVA} + \mathrm{DVA} \]
\end{crdefbox}

\begin{center}
\begin{tikzpicture}[font=\sffamily\footnotesize]
\node[draw=FRMNavy, fill=PaleNavy, line width=0.8pt, rounded corners=1pt,
      text width=3.9cm, align=center, minimum height=1.3cm] (nd) at (0,0)
      {\textbf{No-default value}\\[1pt]$f_{nd}$\\[1pt]\scriptsize from a standard pricing model};
\node[draw=FRMRed, fill=PaleRed, line width=0.8pt, rounded corners=1pt,
      text width=3.6cm, align=center, minimum height=1.3cm] (cva) at (5.4,1.35)
      {\textbf{$-$ CVA}\\[1pt]\scriptsize cost of the \emph{counterparty}\\\scriptsize defaulting};
\node[draw=FRMGreen, fill=PaleGreen, line width=0.8pt, rounded corners=1pt,
      text width=3.6cm, align=center, minimum height=1.3cm] (dva) at (5.4,-1.35)
      {\textbf{$+$ DVA}\\[1pt]\scriptsize cost to the counterparty of \emph{the bank} defaulting};
\node[draw=FRMGold, fill=PaleGold, line width=0.8pt, rounded corners=1pt,
      text width=3.9cm, align=center, minimum height=1.3cm] (ca) at (11.0,0)
      {\textbf{Credit-adjusted value}\\[1pt]$f_{nd}-\mathrm{CVA}+\mathrm{DVA}$};
\draw[-{Stealth[length=5pt]}, FRMRed, line width=0.9pt] (nd) -- (cva);
\draw[-{Stealth[length=5pt]}, FRMGreen, line width=0.9pt] (nd) -- (dva);
\draw[-{Stealth[length=5pt]}, FRMRed, line width=0.9pt] (cva) -- (ca);
\draw[-{Stealth[length=5pt]}, FRMGreen, line width=0.9pt] (dva) -- (ca);
\end{tikzpicture}
\end{center}
\figcap{The two adjustments pull in opposite directions, and the signs are the
examinable part. CVA is subtracted because the counterparty might default on the
bank. DVA is \emph{added} because the bank might default on the counterparty,
which is a benefit to the bank. Note also which party each adjustment is named
for: CVA is about the counterparty's credit, DVA about the bank's own.}

\begin{crfmlbox}[CVA and DVA as discounted sums]
\[ \mathrm{CVA} \;=\; \sum_{i=1}^{N} q_i v_i
   \qquad\qquad
   \mathrm{DVA} \;=\; \sum_{i=1}^{N} q_i^{*} v_i^{*} \]
\vk{The life of the longest outstanding derivative, $T$ years, is divided into
$N$ subintervals. $q_i$ = the probability of a \emph{counterparty} default in
subinterval $i$ and $v_i$ = the present value of the resulting loss; the starred
quantities are the corresponding figures for a \emph{bank} default.}
\end{crfmlbox}

% =====================================================================
\lo{CR-12 c}{Calculate the probability of default using credit spreads.}

This is the hazard rate approximation $\bar{\lambda} = s(T)/(1-\mathrm{RR})$,
worked in full with its bootstrapping at \textbf{CR-9 e}.

% =====================================================================
\lo{CR-12 d}{Describe, compare, and contrast various credit risk mitigants and their role in credit analysis.}

\begin{enumerate}
\item \term{Netting.} Involves offsetting positive and negative cash flows
between two parties, reducing overall exposure. Simplifies risk management and
settlement processes.
\item \term{Collateral agreements.} Involve the exchange of assets, cash or
securities, to mitigate potential losses in the event of default. Marketable
securities are subject to a \term{haircut} to calculate the cash equivalent.
\item \term{Downgrade triggers.} Clauses triggering collateral requirements or
contract termination when a counterparty's credit rating falls below a specified
level.
\end{enumerate}

\begin{crtrapbox}[The limit of a downgrade trigger]
Downgrade triggers \emph{do not} provide protection against a relatively big jump
in a counterparty's credit rating. They work against gradual deterioration, which
is precisely the case in which the counterparty is least likely to fail suddenly.
\end{crtrapbox}

% =====================================================================
\lo{CR-12 e}{Describe the significance of estimating default correlation for credit portfolios and distinguish between reduced form and structural default correlation models.}

Default correlation measures the tendency of multiple entities to default
simultaneously. It is important in the determination of probability distributions
for default losses from a portfolio.

\begin{center}
\footnotesize
\begin{tabular}{L{2.4cm} L{6.0cm} L{6.1cm}}
\toprule
 & \thead{Reduced form models} & \thead{Structural models} \\
\midrule
Mechanism & Assume \emph{stochastic hazard rates} which are correlated with
  macroeconomic factors &
  Based on Merton's model; link defaults to \emph{asset values}. Correlation is
  introduced by assuming the stochastic process followed by the assets of company
  A is correlated with that followed by the assets of company B \\
Correlations produced & \textbf{Low} default correlations &
  \textbf{High} default correlations \\
Practicality & Easy to calculate &
  Time consuming \\
Economic cycle & Reflects economic cycles' impact on default correlations & --- \\
\bottomrule
\end{tabular}
\end{center}
\figcap{The trade-off runs one way down the table: the model that is easy to
compute produces correlations that are too low, and the model that produces
realistic correlations is expensive to run.}

% =====================================================================
\lo{CR-12 f}{Describe the Gaussian copula model for time to default and calculate the probability of default using the one-factor Gaussian copula model.}

To determine the time to default, a Gaussian copula model can be used to quantify
the correlation between the probability distributions of times to default for
multiple companies. Both real-world and risk-neutral default probabilities may be
used.

\begin{crfmlbox}[One-factor Gaussian copula: default probability conditional on the factor]
\[ Q(T \mid F) \;=\; N\!\left(\frac{N^{-1}\bigl[Q(T)\bigr] - \sqrt{\rho}\,F}
   {\sqrt{1-\rho}}\right) \]
\vk{$Q(T)$ = the unconditional cumulative default probability by time $T$; $F$ =
the realization of the common factor; $\rho$ = the copula correlation parameter.
A \emph{low} value of $F$ is a bad state of the world and raises the conditional
default probability.}
\end{crfmlbox}

\begin{crtrapbox}[Why this formula has a minus sign and the WCDR formula has a plus]
The worst case default rate at CR-10 d is written
$N\bigl[(N^{-1}(\mathrm{PD}) + \sqrt{\rho}N^{-1}(X))/\sqrt{1-\rho}\bigr]$, with a
\emph{plus}. This is the same equation. The worst case at confidence $X$ is the
one produced by a sufficiently bad factor realization, that is,
$F = -N^{-1}(X)$, and substituting that into the formula above flips the sign of
the middle term. Both forms are correct; what matters is remembering that $F$
enters \emph{negatively} and the percentile $X$ enters \emph{positively}.
\end{crtrapbox}

% =====================================================================
\lo{CR-12 g}{Describe how to estimate credit VaR using the Gaussian copula and the CreditMetrics approach.}

\begin{crfmlbox}[Credit VaR from the one-factor Gaussian copula]
\[ \text{Credit VaR} \;=\; L\,(1-R)\,\mathrm{WCDR} \]
\vk{$L$ = the size of the loan portfolio; $R$ = the recovery rate; WCDR = the
worst case default rate at the chosen confidence level, computed as at CR-10 d.}
\end{crfmlbox}

\begin{crexbox}[One-factor Gaussian copula on a retail portfolio]
Suppose that a bank has a total of \$100 million of retail exposures. The one-year
probability of default averages 2\% and the recovery rate averages 60\%. The
copula correlation parameter is estimated as 0.1. What is the 99.9\% worst case
default rate, and the one-year 99.9\% credit VaR?
\tcblower
\[ V(0.999, 1) = N\!\left(\frac{N^{-1}(0.02) + \sqrt{0.1}\,N^{-1}(0.999)}
   {\sqrt{1-0.1}}\right) = 0.128 \]
Step by step: $N^{-1}(0.02) = -2.0537$ and $N^{-1}(0.999) = 3.0902$, so the
numerator is $-2.0537 + (0.3162)(3.0902) = -1.0766$, and dividing by
$\sqrt{0.9} = 0.9487$ gives $-1.1348$. Then $N(-1.1348) = 0.1282$.

\smallskip
The one-year 99.9\% credit VaR is therefore
\[ 100 \times 0.128 \times (1 - 0.6) = \$5.13 \text{ million} \]
Note that the unrounded WCDR of 0.12824 is what produces \$5.13 million; using
0.128 exactly would give \$5.12 million.
\end{crexbox}

The \term{CreditMetrics} route to credit VaR is described at CR-10 e: a ratings
transition matrix is built and a Monte Carlo simulation, with a Gaussian copula
governing the joint rating changes, produces the portfolio value distribution
from which VaR is read.

% =====================================================================
